\subsection{How to Tell Whether a Milestone Has Been Met}
\label{sec:10.2.2}
A milestone claimed by the organization that built the system is a different kind of evidence than a milestone checked by somebody else, and the gap between them is not a matter of trust — it is a documented empirical one. Section~\ref{sec:9.1} found that a model's own account of its internal state tracks what its internal computations are actually doing only partially and unreliably, and that the correspondence collapses further depending on how the self-report is elicited. A system's stated confidence that it has reached a milestone is exactly this kind of self-report, and inherits the same unreliability.

A passing score on a fixed test suite fails for a related but different reason. Section~\ref{sec:9.4} covers specification gaming and goal misgeneralization: a method that clears every test in front of it has only been shown to work inside the distribution that test suite covers, and nothing about a passing score distinguishes a system that generalized from one that found whatever pattern the test distribution happened to reward. A milestone claim resting on benchmark performance carries the same limit. What it needs instead is a passing score under conditions built to break it — adversarial, out-of-distribution, or, as section~\ref{sec:9.1}'s account of red-teaming and tamper-resistant safeguards shows for one narrow case, deliberately hostile.

Who runs that harder test matters as much as the test itself. Section~\ref{sec:8.4.2}'s account of self-regulation applies here without much modification: a review conducted or funded by the organization being reviewed answers a structurally different question than one conducted by a party with no stake in the answer, however genuinely the first is intended. METR, a nonprofit that evaluates frontier AI systems' capabilities, runs assessments both in partnership with developers like Anthropic and OpenAI and, separately, on its own after a system has already been released \autocite{metr2026about}. The United Kingdom's AI Security Institute does government-run evaluation of frontier models before deployment and is an independent evaluator, not a rubber stamp — but by its own account it is not a regulator: it cannot compel a company to submit a model for testing, block a release, or intervene once a system is already causing harm \autocite{aisi2025governance}. Independent evaluation and enforcement are not the same capability, and a milestone checked by the first without access to the second is still short of the binding external answer section~\ref{sec:8.4.2} argues self-regulation structurally cannot provide on its own.

Independent does not automatically mean competent to judge. Section~\ref{sec:9.5} covers scalable oversight: once a system's ability in some domain exceeds any available rater's, the rater's verdict, independent or not, is no longer straightforwardly reliable evidence of the property being scored, and no proposed fix for this has been shown to work at the scale it would need to work at. A milestone in a domain where this applies inherits the same limit. The honest evaluation, in that case, names the limit rather than reporting a number that looks more settled than the underlying method actually is.

What an evaluation like this is for is not confirming a milestone but being willing to report that it wasn't met — and, where the milestone turns out to have measured the wrong thing, revising the milestone. A negative result has two different things it could be telling you — that the underlying work fell short, or that the milestone was never a good proxy for the goal it stood in for — and conflating them repeats, one level up, the same self-report failure this section opened with: this time from the system being evaluated to the people who set the terms of its evaluation. 